%%%%%%%%%%%%%%%%%%%%%%%%%%%%%%%%%%%%%%%%%%%%%%
% Journal of Law and Courts target manuscript.
% Official Cambridge/Overleaf small-template source uses cup-journal.
% JLC author instructions: anonymous review, Chicago author-date citations,
% tables/figures near first reference, and replication materials at acceptance.
%%%%%%%%%%%%%%%%%%%%%%%%%%%%%%%%%%%%%%%%%%%%%%
\documentclass[
  journal=small,
  manuscript=research-article,
  year=2026
]{cup-journal}

\usepackage{amsmath}
\usepackage[nopatch]{microtype}
\usepackage{array}
\usepackage{booktabs}
\usepackage{graphicx}
\usepackage{longtable}
\usepackage[authoryear,round]{natbib}
\usepackage{url}
\usepackage{xcolor}
\usepackage[hidelinks]{hyperref}

\makeatletter
\newcolumntype{L}[1]{>{\raggedright\arraybackslash}p{#1}}
\newenvironment{inlinefigureblock}{%
  \par\smallskip
  \begingroup
  \centering
}{%
  \par\smallskip
  \endgroup
}

\newcommand{\inlinefigurecaption}[2]{%
  \refstepcounter{figure}%
  \label{#2}%
  \vskip\abovecaptionskip
  \@makecaption{\figurename~\thefigure}{#1}%
  \vskip\belowcaptionskip
}
\makeatother

\title{Stress-Testing Constitutional Review Design Bundles: A Comparative Simulation Framework}

\author{Anonymous Author}
\affiliation{Affiliation withheld for review}
\date{}

\keywords{constitutional review, judicial politics, comparative courts, institutional design, simulation}

\begin{document}
\maketitle

\begin{abstract}
This article introduces a simulation framework for stress-testing constitutional-review design bundles under shared generated dockets. It compares how appointment, tenure, emergency procedure, screening, legislative response, compliance, and cost shape rights protection, democratic responsiveness, and veto-relocation risk. A frozen test of 66 Supreme Court Database disputes from the 2025 term finds no clear conditional Brier-score advantage for a training-only empirical docket profile over the fixed generator. A paired study of four forecast profiles, thirteen contexts, and eight base or modified designs then qualifies three institutional claims. Reasoned emergency review lowers modeled irregularity, but the shared review-time input cannot identify a procedural-duration effect. Council upstream-cost differences remain below the practical threshold. Disabling legislative response weakens some weak-form outcomes but leaves other modeled advantages intact. These are findings about measurement and model dependence, not identified reform effects. Historical anchoring concerns docket composition only; source-range checks, conditional Monte Carlo intervals, structural sensitivity, and illustrative normative preferences remain separate evidence layers. The contribution is an auditable account of which conclusions the model supports, weakens, or cannot measure, rather than a ranking of real courts.
\end{abstract}

\section{Introduction}

Constitutional review systems face a recurring design problem: courts need enough independence to enforce constitutional limits, but enough accountability and procedural discipline to avoid becoming unreviewable political actors. The classic countermajoritarian framing emphasizes the difficulty of judicial invalidation in a democracy \citep{bickel1962}, and court-skeptical accounts press the democratic case for limiting judicial supremacy \citep{tushnet1999}. Comparative constitutional scholarship shows that review can be organized through many institutional forms rather than a single supreme-court template \citep{stoneSweet2000,ginsburg2003}. Law-and-courts scholarship also treats courts as institutions whose behavior depends on selection rules, legal hierarchy, compliance expectations, and political context \citep{epsteinKnight1998,helmke2005,vanberg2005,staton2010}.

The article's main contribution is methodological. It treats constitutional review as a bundled institutional system rather than as a single yes-or-no choice about judicial review. The framework compares how design choices shift authority across courts, councils, screeners, emergency procedures, legislatures, and enforcement channels while holding the generated worlds and dockets constant. The intervention is therefore in the debate over judicial review and democratic constitutionalism: instead of asking only whether courts have too much or too little power, the article asks where final blocking authority moves when reformers change review institutions.

That shift is the paper's central concern. Veto-relocation risk arises when a reform reduces one visible form of judicial finality while moving blocking power into a less visible or less accountable stage. The problem is not that all relocation is bad. Some relocation may improve rights screening, reason-giving, or legislative response. The risk appears when the new veto point is harder to contest, less transparent, less reasoned, or more dependent on unequal access.

The framework evaluates this risk across stylized institutional archetypes: U.S.-style discretionary apex review, Kelsenian constitutional-court review, council-style pre-enactment and abstract review, weak-form parliamentary review, supranational treaty review, and dual-court or cross-checking review. Those archetypes discipline the set of institutional designs and make institutional tradeoffs visible.

A secondary comparison imports legislative-output profiles as docket stress signals, preserving subject matter where possible and testing whether results change under high-capture, high-volatility, and low-mandate docket conditions.

The central question is when constitutional review improves democratic constitutionalism rather than merely relocating veto power. Baseline comparisons suggest possible emergency-process, screening-burden, and legislative-response tradeoffs. The historical robustness study tests those suggestions and narrows them: duration is not a design-specific outcome, substantial upstream displacement is not established by the council-structure intervention, and response dependence explains only part of weak-form advantages. Section~\ref{sec:infrastructure} defines constitutional review as institutional infrastructure. Section~\ref{sec:framework} describes the simulation framework. Section~\ref{sec:experiments} summarizes the original comparative stress tests. Section~\ref{sec:historical} reports the frozen historical evaluation and paired robustness study. Section~\ref{sec:diagnostics} reports source-range checks, uncertainty, and replication limits. Sections~\ref{sec:implications} and \ref{sec:limits} address institutional interpretation and remaining limitations.

\section{Constitutional Review as Institutional Infrastructure}
\label{sec:infrastructure}

A constitutional review system is more than an apex court. It includes appointment and replacement rules, lower-court pathways, docket filters, emergency relief procedures, merits review, voting thresholds, opinion coalitions, recusal practices, legislative override rules, and enforcement channels. Treating these features as separate design variables makes sense analytically, but reform proposals usually bundle them. The simulation therefore models review as an infrastructure problem: a constitutional system routes disputes through a sequence of institutions, and the consequences depend on how those routing rules interact.

The model components are tied to distinct literatures rather than to a single court-reform theory. Political-insurance and juristocracy accounts motivate treating constitutional review adoption and design as political-institutional choices \citep{hirschl2004,ginsburgVersteeg2014}. Weak-form and Commonwealth-model scholarship motivates declaration, override, and legislative-response mechanisms \citep{gardbaum2013}. Access-to-justice and support-structure work motivates the case-selection-access and public-defender measures \citep{epp1998}. Comparative judicial-role scholarship motivates the attention to compliance, implementation, and political consequences \citep{kapiszewski2013}. Emergency-docket scholarship motivates separating interim relief, reasons, vote disclosure, public disagreement, and merits follow-up \citep{baude2015,vladeck2023}.

Table~\ref{tab:design-space} narrows the design space used in this paper. The simulator contains additional scenario-library variants, but the manuscript focuses on four mechanism families and one archetype family so that the reader can tell what is actually being compared.

\begin{longtable}{L{0.24\linewidth}L{0.30\linewidth}L{0.34\linewidth}}
\caption{Scenario groups used in the reported stress tests.}
\label{tab:design-space}\\
\toprule
Scenario group & Examples in reported comparison & Primary tradeoff \\
\midrule
\endfirsthead
\toprule
Scenario group & Examples in reported comparison & Primary tradeoff \\
\midrule
\endhead
Court-architecture variants & Current-style discretionary apex court; 18-year terms; commission appointment; supermajority invalidation; strict recusal; panel/en banc routing; dual cross-checking courts; constitutional council; legislative override; retention accountability; hybrid independence/accountability. & Whether selection, size, threshold, routing, and accountability reforms change rights, legitimacy, conflict, cost, and veto-relocation risk when the docket is held fixed. \\
Emergency-review procedure & Current-style emergency orders, strict recusal, reasoned emergency panels, cross-checking or council review, and hybrid review. & Whether urgent rights protection can be separated from opaque or merits-displacing emergency intervention. \\
Dialogic and weak-form mechanisms & Weak-form review, suspended declarations, override clauses, and mandatory legislative response cycles. & Whether review generates credible democratic response rather than final judicial veto or unresolved legislative drift. \\
Front-end and access mechanisms & Pre-enactment review, abstract review, ombudsman-triggered review, constitutional public defenders, and rights-impact statements. & Whether access and early screening improve rights-sensitive review or hide veto power in upstream institutions. \\
Stylized archetype presets & U.S.-style discretionary apex review, Kelsenian constitutional-court review, council-style pre-enactment or abstract review, weak-form parliamentary review, supranational treaty review, and dual-court/cross-checking review. & Source-range comparison and scenario discipline. \\
Feasibility constraints & Legal-transplant feasibility, implementation cost, political-culture sensitivity, implementation capacity, and case-selection access. & Whether a reform remains plausible once cost, administrative load, legal culture, compliance, and unequal access are visible. \\
\bottomrule
\end{longtable}

The central design question is not whether one feature is good or bad in isolation. It is whether a bundle produces an acceptable balance among legal stability, rights protection, democratic responsiveness, legitimacy, conflict management, administrative burden, and enforceability.

The paper distinguishes court architecture scenarios from mechanism scenarios. Architecture scenarios approximate institutional bundles; mechanism scenarios isolate specific reform devices. Source-range checks compare selected outputs with documented external ranges and identify calibration gaps.

\section{Simulation Framework}
\label{sec:framework}

The simulator uses named institutional designs, randomized worlds, shared dockets, repeated simulation comparisons, directional measures, structured outputs, and replication records. Each case is generated from five input families: legal stakes, procedural posture, political context, access capacity, and institutional pressure. The full variable list and equations are reported in the supplementary appendix.

Figure~\ref{fig:model-pipeline} summarizes the model pipeline. Legislative-output profiles enter only as structured stress signals that alter the generated constitutional-review docket.

\begin{figure}[hbt!]
\centering
% Auto-generated by paper/scripts/generate_figures.py
\begingroup
\setlength{\unitlength}{1mm}
\begin{picture}(128,66)
\scriptsize
\put(2,45){\framebox(22,12){\shortstack[c]{Scenario\\catalog}}}
\put(28,45){\framebox(22,12){\shortstack[c]{World and\\case generation}}}
\put(54,45){\framebox(22,12){\shortstack[c]{Institutional\\routing}}}
\put(80,45){\framebox(22,12){\shortstack[c]{Decision\\process}}}
\put(106,45){\framebox(20,12){\shortstack[c]{Diagnostics\\and reports}}}
\put(24,51){\vector(1,0){4}}
\put(50,51){\vector(1,0){4}}
\put(76,51){\vector(1,0){4}}
\put(102,51){\vector(1,0){4}}
\put(28,26){\framebox(22,10){\shortstack[c]{Legislative\\outputs}}}
\put(39,36){\vector(0,1){9}}
\put(54,28){\framebox(22,10){\shortstack[c]{Intake and\\access filters}}}
\put(80,28){\framebox(22,10){\shortstack[c]{Emergency\\and merits}}}
\put(106,28){\framebox(20,10){\shortstack[c]{Compliance\\reaction cost}}}
\put(50,33){\vector(1,0){4}}
\put(76,33){\vector(1,0){4}}
\put(102,33){\vector(1,0){4}}
\put(116,28){\vector(0,-1){8}}
\put(64,12){\framebox(38,8){\shortstack[c]{Period turnover and\\political-culture state}}}
\put(64,16){\vector(-1,0){40}}
\put(83,20){\vector(0,1){8}}
\put(13,38){\makebox(0,0){\shortstack[c]{appointment, size, tenure,\\thresholds, councils, overrides}}}
\put(115,19){\makebox(0,0){\shortstack[c]{direct outputs, composites,\\uncertainty, provenance}}}
\put(64,5){\makebox(0,0){\shortstack[c]{Case generation, routing, decisions, compliance, reaction, and scoring are separate audit points.}}}
\end{picture}
\endgroup

\caption{Simulation pipeline and secondary legislative-output stress bridge.}
\label{fig:model-pipeline}
\end{figure}

Each scenario creates justices or review actors from institutional rules such as selection, size, tenure, removal, transparency, review archetype, timing, cost profile, and review mechanism. The review process then models intake, emergency handling, recusals, panel or en banc review, voting thresholds, cross-checking, council screening, weak-form or suspended declarations, legislative response, overrides, access mechanisms, and opinion fragmentation. The comparative scenarios use fixed weights rather than fitted parameters; the named source-range presets add bounded calibration hooks for denominator conventions, emergency procedure, deferred remedies, and weak-form response checks so that the external audit compares like to like.

Interpretation rule. Court-architecture scenarios are stylized institutional comparators. Mechanism scenarios are counterfactual design tests. Source-range checks compare selected model outputs with documented external ranges. The comparison is therefore designed to isolate tradeoffs inside the specified model, not to forecast named courts.

The weights have different epistemic status. Some inputs are institutional mappings, such as court size, route type, and whether a design uses weak-form declarations or merits follow-up. Some are source-observed benchmark ranges. Some are literature-informed directional priors, such as the expectation that political conflict reduces compliance and reason-giving improves process transparency. Some are design-prior weights used to make mechanisms comparable under shared worlds. The analysis below treats source checks, mechanism movement, and normative score construction as separate claims.

The democratic-constitutionalism composite combines rights protection, democratic responsiveness, legality, legitimacy, compliance, low conflict, low veto-relocation risk, access, implementation capacity, transplant feasibility, and political-culture robustness because those are the values most directly implicated by the design bundles modeled here. Different constitutional theories would weight those values differently. The experiments therefore report direct outputs first, use the composite score as a display aid, and add alternative normative profiles.

\begin{table}[hbt!]
\centering
\caption{Model-input transparency and parameter status.}
\label{tab:model-transparency}
\scriptsize
\setlength{\tabcolsep}{3pt}
\begin{tabular}{@{}L{0.21\linewidth}L{0.25\linewidth}L{0.23\linewidth}L{0.21\linewidth}@{}}
\toprule
Input family & Scale or construction & Role in model & Status in this paper \\
\midrule
Generated case facts & Bounded $[0,1]$ draws with correlations among rights threat, urgency, salience, ambiguity, public support, lower-court conflict, and policy domain. & Creates the shared docket faced by each scenario. & Design prior; stress-tested through comparison cases and legislative-output imports. \\
World context & Bounded controls for appointment polarization, rights-threat rate, emergency pressure, legislative conflict, public trust, partisan pressure, fragmentation, government control, electoral time pressure, civil society, implementation capacity, and legal-tradition fit. & Shapes docket stress, replacement dynamics, compliance, and response. & Stress-case parameter; country-context rows support source-range comparisons. \\
Institutional rules & Scenario-coded court size, appointment method, tenure, removal, recusal, routing, emergency procedure, voting threshold, council/cross-checking, override, and weak-form mechanisms. & Determines intake, voting, emergency handling, response, cost, and veto-relocation channels. & Institutional mapping or mechanism flag. \\
Model weights & Stated coefficients for intake, emergency relief, merits invalidation, compliance, response, conflict, legitimacy, cost, and composites. & Converts case facts and rules into outputs. & Explicit design-prior weights; not fitted coefficients. Alternative profile and stress-case results test sensitivity. \\
External source ranges & Source-backed rows with source URL, denominator, construction note, and direct simulator analogue; contextual rows are retained but excluded from fit counts. & Finds missing mechanisms and denominator problems. & Source-range comparison. \\
\bottomrule
\end{tabular}
\end{table}

The main model equations are intentionally simple enough to audit; the supplementary appendix states the fuller notation, composite-score equations, and review-path construction. Review propensity is a weighted score:
\[
\begin{aligned}
R={}&A_s+.26T_r+.22S_c+.16L_c+.22C_p+.12I_c+.18E\\
&+.16O+.13D+.10P+.12B+.09L_a+.12P_i+.05G_r-.05U,
\end{aligned}
\]
where $A_s$ is structural access, $T_r$ rights threat, $S_c$ constitutional salience, $L_c$ lower-court conflict, $C_p$ certiorari or intake pressure, $I_c$ intercourt conflict, $E$ emergency status, $O$ ombudsman trigger, $D$ public-defender participation, $P$ pre-enactment review, $B$ abstract review, $L_a$ litigant capacity, $P_i$ public-interest support, $G_r$ government repeat-player advantage, and $U$ legal ambiguity. The case is reviewed when $R$ multiplied by jurisdictional access plus random docket noise exceeds the review threshold. Government repeat-player advantage can raise case salience and review probability while still lowering case-selection access and democratic-constitutionalism scores when it reflects asymmetric litigating capacity.

The review vote uses justice-level rights sensitivity, partisan pressure, legal ambiguity, lower-court signal, mandate restraint, and mechanism adjustments. Mechanism adjustments are deliberately directional rather than fitted: rights-impact statements slightly reduce unnecessary invalidation while increasing attention to high-rights-threat cases; ombudsman and public-defender participation increase access and rights salience; pre-enactment and abstract review increase salience screening; weak-form review lowers the threshold for a declaration without automatically invalidating the law.

Emergency orders are modeled separately from merits invalidation. The simulator tracks interim relief, reason-giving, vote disclosure, public disagreement, government applicants, government wins, and merits follow-up. This distinction matters because an institution can be restrained on the merits while still aggressive on emergency relief, or transparent in merits decisions while opaque on emergency orders.

Figure~\ref{fig:emergency-docket} illustrates why that separation matters. The current-style scenario generates high emergency relief and high emergency-process irregularity, while reason-giving and cross-checking designs can keep emergency relief available while sharply changing transparency and follow-up behavior.

\begin{figure}[hbt!]
\centering
% Auto-generated by paper/scripts/generate_figures.py
\begingroup
\setlength{\unitlength}{1mm}
\begin{picture}(128,72)
\scriptsize
\put(36.0,10){\color{black!15}\line(0,1){50}}
\put(36.0,6){\makebox(0,0){0}}
\put(55.5,10){\color{black!15}\line(0,1){50}}
\put(55.5,6){\makebox(0,0){0.25}}
\put(75.0,10){\color{black!15}\line(0,1){50}}
\put(75.0,6){\makebox(0,0){0.5}}
\put(94.5,10){\color{black!15}\line(0,1){50}}
\put(94.5,6){\makebox(0,0){0.75}}
\put(114.0,10){\color{black!15}\line(0,1){50}}
\put(114.0,6){\makebox(0,0){1}}
\put(36.0,10){\line(1,0){78.0}}
\put(36.0,10){\line(0,1){50}}
\put(36,64){\color{black}\rule{6mm}{1.2mm}}
\put(44,64){\makebox(0,0)[l]{shadow-docket abuse}}
\put(76,64){\color{black!55}\rule{6mm}{1.2mm}}
\put(84,64){\makebox(0,0)[l]{reason-giving}}
\put(36,60.5){\color{black!25}\rule{6mm}{1.2mm}}
\put(44,60.5){\makebox(0,0)[l]{merits follow-up}}
\put(32,56.0){\makebox(0,0)[r]{\color{red}CUR}}
\put(36.0,57.8){\color{black}\rule{28.5mm}{1.2mm}}
\put(36.0,56.0){\color{black!55}\rule{2.0mm}{1.2mm}}
\put(36.0,54.2){\color{black!25}\rule{0.0mm}{1.2mm}}
\put(32,48.0){\makebox(0,0)[r]{\color{black}REC}}
\put(36.0,49.8){\color{black}\rule{4.8mm}{1.2mm}}
\put(36.0,48.0){\color{black!55}\rule{54.4mm}{1.2mm}}
\put(36.0,46.2){\color{black!25}\rule{22.5mm}{1.2mm}}
\put(32,40.0){\makebox(0,0)[r]{\color{black}EMR}}
\put(36.0,41.8){\color{black}\rule{0.2mm}{1.2mm}}
\put(36.0,40.0){\color{black!55}\rule{67.7mm}{1.2mm}}
\put(36.0,38.2){\color{black!25}\rule{78.0mm}{1.2mm}}
\put(32,32.0){\makebox(0,0)[r]{\color{black}DUAL}}
\put(36.0,33.8){\color{black}\rule{0.2mm}{1.2mm}}
\put(36.0,32.0){\color{black!55}\rule{69.9mm}{1.2mm}}
\put(36.0,30.2){\color{black!25}\rule{78.0mm}{1.2mm}}
\put(32,24.0){\makebox(0,0)[r]{\color{black}COUNC}}
\put(36.0,25.8){\color{black}\rule{4.6mm}{1.2mm}}
\put(36.0,24.0){\color{black!55}\rule{58.3mm}{1.2mm}}
\put(36.0,22.2){\color{black!25}\rule{22.9mm}{1.2mm}}
\put(32,16.0){\makebox(0,0)[r]{\color{black}HYB}}
\put(36.0,17.8){\color{black}\rule{0.2mm}{1.2mm}}
\put(36.0,16.0){\color{black!55}\rule{68.2mm}{1.2mm}}
\put(36.0,14.2){\color{black!25}\rule{78.0mm}{1.2mm}}
\put(75,2){\makebox(0,0){rate}}
\end{picture}
\endgroup

\caption{Emergency-docket outputs in the baseline comparison. Higher reason-giving and merits follow-up are generally desirable; lower emergency-process irregularity is generally desirable.}
\label{fig:emergency-docket}
\end{figure}

A simple toy case clarifies the state transition. Suppose the docket contains a high-rights-burden election dispute with high urgency, weak vehicle quality, substantial lower-court conflict, and strong executive pressure. Under a fast emergency docket, the case can receive interim relief before merits review, producing rights-claimant success in the short run but also raising emergency-process irregularity when reasons, vote disclosure, or merits follow-up are absent. Under a written-reasons design, the same emergency relief is more transparent, so the direct opacity component falls, but the model still asks whether later merits review, compliance, and conflict improve. Under a no-relief-without-merits connection, the case is forced toward expedited merits handling; modeled irregularity can fall while administrative burden or the probability of a denied urgent claimant changes. Actual procedural duration is not identified: the reported average review time is a shared case input, and the separate delay-cost index is not elapsed time. This example motivates examining downstream compliance, conflict, cost, and rights-claimant success without treating direct rule effects as independent validation.

A dialogic-review vignette exposes the same problem in a different form. Suppose a high-mandate legislature enacts a rights-burdening criminal-procedure statute. A strong-form court can invalidate the statute immediately, concentrating final veto authority in the judiciary. A weak-form court can issue a declaration while giving the legislature a response period. Whether that response occurs is a substantive question, not something a declaration alone establishes. If the legislature ignores the declaration or issues a symbolic reply, authority can move from a court order to legislative inaction. The response-off study below tests this dependency while leaving the design's other terms intact. Rights-impact statements pose a related question about whether front-end screening remains transparent and contestable.

\begin{table}[hbt!]
\centering
\caption{Mechanical implications and downstream tests.}
\label{tab:mechanical-diagnostic}
\scriptsize
\setlength{\tabcolsep}{3pt}
\begin{tabular}{@{}L{0.24\linewidth}L{0.24\linewidth}L{0.42\linewidth}@{}}
\toprule
Modeled pattern & Mechanically implied? & What would make it substantively interesting \\
\midrule
Written emergency reasons reduce the opacity portion of emergency irregularity & Partly yes & Whether reasons also change merits follow-up, compliance, conflict, public-legitimacy proxy, or administrative cost. \\
No-relief-without-merits rules reduce merits-displacing emergency action & Partly yes & Whether urgent rights-claimant success is preserved and whether the downstream emergency effect falls without unacceptable delay. \\
Appointment or size reforms move emergency outputs less directly & No direct coding implication & Whether composition changes alter emergency outcomes after shared dockets, legal ambiguity, replacement timing, and public trust are held constant. \\
Weak-form or response-cycle mechanisms improve democratic constitutionalism & Partly score- and design-dependent & Whether changing actual response opportunities moves downstream outcomes after other design terms are held fixed. \\
\bottomrule
\end{tabular}
\end{table}

Runs are partitioned into review periods so term limits and removal or retention rules can change court composition over time. Decision outcomes also update an institutional reaction state: public trust, legislature-court conflict, court-curbing pressure, override pressure, amendment pressure, and compliance norms feed later cases in the same run. Enforcement is split into executive implementation, agency nonacquiescence, legislative reenactment, and local-government compliance. World assumptions include party fragmentation, governing coalition control, electoral time pressure, civil-society capacity, implementation capacity, and legal-tradition compatibility.

The main outcome metrics are listed in Table~\ref{tab:metrics}. The table separates direct procedural outputs, direct outcome outputs, constructed indices, and display-only reading aids. That distinction matters. Emergency reason-giving, vote disclosure, merits follow-up, weak-form declarations, legislative responses, and public-defender participation are direct model events. Rights-claimant success and invalidation are direct outcomes, but they are not identical to morally correct rights protection; a successful claimant can be strategically selected, anti-regulatory, or in tension with other rights holders. Legitimacy is therefore labeled as a modeled public-legitimacy proxy rather than an opinion-survey measure, and the legacy \texttt{shadowDocketAbuse} report column is interpreted as emergency-process irregularity rather than a legal conclusion that a specific order is abusive. Constitutional conflict is likewise a composite index covering legal, political, federalism, and implementation conflict.

The constitutional-design measures make three tradeoffs explicit. Veto-relocation risk rises when a design shifts blocking power into courts, councils, cross-checks, or upstream admissibility screens without a rights-threat justification or legislative response. Legal-transplant feasibility combines scenario transplant score, political-culture fit, trust, inverse institutional demand, legal-tradition compatibility, implementation capacity, civil-society capacity, and inverse political mismatch. Political-culture sensitivity increases when a mechanism depends heavily on cooperative legislative response, high public trust, low polarization, low fragmentation, or strong compliance norms. Democratic constitutionalism combines the values most central to the paper's question:
\[
\begin{aligned}
D_c={}&.11S+.18R_p+.12L+.15D_r+.09C+.08(1-K)+.07(1-V)+.05F\\
&+.03(1-P_c)+.03Q+.03A_c+.03(1-G_r)+.03I_m,
\end{aligned}
\]
where $S$ is legal stability, $R_p$ rights protection, $L$ the modeled public-legitimacy proxy, $D_r$ democratic responsiveness, $C$ compliance, $K$ the constitutional-conflict index, $V$ veto-relocation risk, $F$ transplant feasibility, $P_c$ political-culture sensitivity, $Q$ legislative-response credibility, $A_c$ case-selection access, $G_r$ government repeat-player advantage, and $I_m$ implementation capacity.

\begin{table}[hbt!]
\centering
\caption{Core metric families.}
\label{tab:metrics}
\begin{tabular}{@{}L{0.23\linewidth}L{0.50\linewidth}L{0.13\linewidth}@{}}
\toprule
Metric family & Examples & Direction \\
\midrule
Legal performance & legal stability, rights protection, reversal rate, merits invalidation rate & mixed \\
Political alignment & partisan alignment, democratic responsiveness, independence/accountability balance & mixed \\
Democratic constitutionalism & democratic constitutionalism, legislative-response credibility, case-selection access, government repeat-player advantage, implementation capacity, veto-relocation risk, transplant feasibility, political-culture sensitivity & mixed \\
Noncourt mechanisms & weak-form declarations, suspended declarations, legislative response cycles, rights-impact statements, ombudsman triggers, public-defender participation, pre-enactment review, abstract review & reported \\
Supranational route mix & preliminary references, appeal routes, direct actions & reported \\
Emergency docket & emergency relief, reason-giving, vote disclosure, public disagreement, merits follow-up & mixed \\
Legitimacy and conflict & modeled public-legitimacy proxy, public trust input, constitutional-conflict index, court-curbing pressure, amendment pressure & mixed \\
Pipeline and intake & review rate, intake acceptance, case-selection access, litigant capacity, public-interest support, lower-court depth, state-federal tension, intercourt conflict & reported \\
Coalitions and recusals & concurrence fragmentation, dissent intensity, recusal rate, en banc rate, cross-check rate & mixed \\
Compliance & compliance, defiance, workarounds, executive implementation, agency nonacquiescence, reenactment, local compliance & mixed \\
Institutional cost & administrative load, direct court cost, upstream screening cost, capacity strain, delay, implementation complexity & lower \\
\bottomrule
\end{tabular}
\end{table}

\subsection{Directional score construction}

The directional score is a display index. It is useful for locating clusters of designs, but it is not a welfare function, an empirical estimate, or a basis for fine ordinal ranking. Let $H$ be the set of higher-is-better metrics and $L$ the set of lower-is-better metrics included in the display score. The score is
\[
S=\frac{1}{35}\left(\sum_{m\in H}m+\sum_{m\in L}(1-m)\right).
\]
The higher-is-better set includes legal stability, rights protection, modeled public-legitimacy proxy, democratic responsiveness, legislative-response credibility, timely legislative response, case-selection access, implementation capacity, democratic constitutionalism, independence/accountability balance, compliance, executive implementation, local-government compliance, emergency reason-giving, emergency vote disclosure, legal-transplant feasibility, legislative response, and rights-impact statements. The lower-is-better set includes partisan alignment, emergency-process irregularity, reversal rate, emergency relief rate, merits invalidation rate, constitutional-conflict index, government repeat-player advantage, veto-relocation risk, political-culture sensitivity, defiance, workarounds, agency nonacquiescence, legislative reenactment, emergency public disagreement, legislative-response delay, administrative load, and total institutional cost. Reported metrics such as review rate, recusal rate, override rate, weak-form declaration rate, abstract-review rate, lower-court depth, replacement pressure, and public trust remain visible in reports but do not enter $S$.

Designs with overlapping bootstrap intervals or score differences below about one percentage point should be treated as tied unless an individual mechanism output explains a stable difference. The main results therefore privilege multi-objective comparisons--emergency irregularity, rights-claimant success, democratic constitutionalism, public-legitimacy proxy, veto-relocation risk, transplant feasibility, political-culture sensitivity, and cost--over score order.

\section{Comparative Experiments}
\label{sec:experiments}

The main simulation comparison evaluates court-architecture scenarios and mechanism scenarios across five stress cases: baseline, polarized appointments, rights-threat surge, emergency-docket stress, and low-trust conflict. A sensitivity comparison reruns the institutional designs under high and low assumptions for emergency pressure, appointment polarization, rights-threat rate, public trust, and legislature-court conflict.

\subsection{Reading the outputs}

The results separate direct implementation checks from downstream propagation effects. Direct checks include reason-giving, merits follow-up, mechanism activation, weak-form declarations, suspended declarations, and formal legislative-response triggers. They show whether the mechanism is operating as coded. Downstream outputs include rights-claimant success, compliance, constitutional conflict, modeled public legitimacy, institutional cost, and veto-relocation risk. These are still consequences of assigned model rules, not empirical outcome estimates. Average review time is a shared input rather than a downstream duration measure.

Table~\ref{tab:baseline-results} reports the baseline estimates that anchor the graphical summaries. The table puts direct and constructed outputs before the display score so that emergency-process irregularity, rights-claimant success, modeled legitimacy, veto relocation, feasibility, and cost remain visible.

% Auto-generated by paper/scripts/generate_tables.py
\begin{table}[hbt!]
\centering
\caption{Court-architecture stress-test results.}
\label{tab:baseline-results}
\scriptsize
\setlength{\tabcolsep}{2pt}
\begin{tabular}{@{}lL{0.26\linewidth}rrrrrrr@{}}
\toprule
Code & Architecture & Dem. const. & Rights & Emerg. irr. $\downarrow$ & Legit. proxy & Veto $\downarrow$ & Cost $\downarrow$ & Score aid \\
\midrule
CUR & Current-style discretionary apex court & 0.571 & 0.687 & 0.366 & 0.582 & 0.529 & 0.450 & 0.525 \\
18Y & Eighteen-year terms & 0.578 & 0.695 & 0.359 & 0.592 & 0.520 & 0.455 & 0.535 \\
15J & Fifteen-justice commission & 0.643 & 0.739 & 0.062 & 0.708 & 0.436 & 0.470 & 0.615 \\
SUP & Supermajority invalidation rule & 0.638 & 0.736 & 0.064 & 0.684 & 0.446 & 0.457 & 0.604 \\
REC & Strict recusal court & 0.648 & 0.739 & 0.063 & 0.726 & 0.434 & 0.476 & 0.622 \\
EMR & Reasoned emergency review & 0.649 & 0.740 & 0.003 & 0.740 & 0.446 & 0.460 & 0.625 \\
PAN & Panel and en banc routing & 0.643 & 0.739 & 0.062 & 0.720 & 0.433 & 0.511 & 0.609 \\
DUAL & Dual cross-checking courts & 0.650 & 0.747 & 0.003 & 0.751 & 0.596 & 0.589 & 0.628 \\
COUNC & Constitutional council & 0.652 & 0.741 & 0.058 & 0.753 & 0.662 & 0.533 & 0.632 \\
OVR & Legislative override court & 0.637 & 0.733 & 0.064 & 0.700 & 0.435 & 0.466 & 0.592 \\
RET & Retention-accountability court & 0.656 & 0.729 & 0.057 & 0.723 & 0.407 & 0.482 & 0.622 \\
HYB & Independence-accountability hybrid & 0.663 & 0.746 & 0.002 & 0.754 & 0.433 & 0.542 & 0.640 \\
\bottomrule
\end{tabular}
\begin{minipage}{0.96\linewidth}
\footnotesize Notes: Values are baseline campaign estimates for court-architecture scenarios only; mechanism scenarios are separated in Table~\ref{tab:mechanism-results}. Lower emergency irregularity indicates less opaque or merits-displacing emergency intervention. The legitimacy column is a modeled public-legitimacy proxy. The score aid is the equally weighted directional display score; its 95\% run-block bootstrap interval is reported in the replication files.
\end{minipage}
\end{table}


Table~\ref{tab:mechanism-results} reports mechanism activation and the downstream propagation columns in the same row. This format is deliberate: for example, the important question is not merely whether a weak-form scenario triggers declarations, but whether those declarations also improve response credibility, rights protection, compliance, conflict, cost, or veto-relocation risk. The front-end rows are not pure single-mechanism toggles. In the current implementation, rights-impact statements, pre-enactment review, and abstract review are routing bundles, so a row labeled by one front-end mechanism can also show realized activation of the others.

% Auto-generated by paper/scripts/generate_tables.py
\begin{table}[hbt!]
\centering
\caption{Mechanism activation and downstream propagation effects.}
\label{tab:mechanism-results}
\scriptsize
\setlength{\tabcolsep}{2pt}
\begin{tabular}{@{}lL{0.18\linewidth}L{0.22\linewidth}rrrrrrr@{}}
\toprule
Code & Mechanism & Direct activation & Rights & Compliance & Conflict $\downarrow$ & Legit. & Cost $\downarrow$ & Veto $\downarrow$ & Dem. const. \\
\midrule
WFR & Weak-form review & Weak 0.357; Resp. 0.376 & 0.718 & 0.802 & 0.208 & 0.754 & 0.359 & 0.177 & 0.709 \\
SUS & Suspended declaration & Susp. 0.442; Resp. 0.451 & 0.724 & 0.596 & 0.246 & 0.761 & 0.383 & 0.207 & 0.680 \\
OVC & Override clause & Override 0.000; Resp. 0.261 & 0.738 & 0.605 & 0.221 & 0.737 & 0.425 & 0.355 & 0.655 \\
PRE & Pre-enactment review & RIS 0.381; Pre 0.833; Abs 0.872 & 0.766 & 0.683 & 0.234 & 0.755 & 0.437 & 0.511 & 0.667 \\
ABS & Abstract review tribunal & RIS 0.381; Pre 0.833; Abs 0.872 & 0.770 & 0.603 & 0.234 & 0.747 & 0.436 & 0.312 & 0.665 \\
OMB & Ombudsman-triggered review & Omb. 0.263 & 0.736 & 0.533 & 0.252 & 0.729 & 0.516 & 0.525 & 0.632 \\
DEF & Constitutional public defender & Def. 0.306 & 0.773 & 0.537 & 0.240 & 0.741 & 0.518 & 0.514 & 0.646 \\
RIS & Rights-impact statement & RIS 1.000; Pre 0.833; Abs 0.872 & 0.805 & 0.720 & 0.231 & 0.773 & 0.437 & 0.502 & 0.691 \\
MLR & Mandatory legislative response & Resp. 0.390 & 0.717 & 0.665 & 0.197 & 0.782 & 0.394 & 0.224 & 0.680 \\
\bottomrule
\end{tabular}
\begin{minipage}{0.96\linewidth}
\footnotesize Notes: Direct activation reports implementation-facing outputs; the remaining columns report downstream propagation effects. Front-end rows are bundle/routing scenarios rather than isolated toggles, so RIS, Pre, and Abs may all appear when a design routes review through multiple upstream screens. Lower conflict, cost, and veto-relocation risk are better; higher rights, compliance, legitimacy, and democratic constitutionalism are better.
\end{minipage}
\end{table}


Table~\ref{tab:normative-profiles} makes the score sensitivity problem explicit. It uses the same baseline outputs but changes the display weights to match five plausible constitutional priorities. The preferred architecture or mechanism changes when the normative theory changes.

% Auto-generated by paper/scripts/generate_tables.py
\begin{table}[hbt!]
\centering
\caption{Alternative normative-profile readings of the same baseline outputs.}
\label{tab:normative-profiles}
\scriptsize
\setlength{\tabcolsep}{3pt}
\begin{tabular}{@{}L{0.18\linewidth}L{0.24\linewidth}L{0.20\linewidth}rL{0.20\linewidth}@{}}
\toprule
Profile & Emphasis & Leading architecture & Score & Leading mechanism \\
\midrule
Rights-centered & rights, access, compliance, implementation & HYB (Independence-accountability hybrid) & 0.622 & WFR (Weak-form review) \\
Parliamentary-dialogue & responsiveness, response credibility, low conflict & RET (Retention-accountability court) & 0.492 & SUS (Suspended declaration) \\
Legal-stability & stability, compliance, low conflict & COUNC (Constitutional council) & 0.759 & WFR (Weak-form review) \\
Cost-minimizing & low cost, low load, implementation capacity & SUP (Supermajority invalidation rule) & 0.566 & WFR (Weak-form review) \\
Anti-veto-relocation & visible access with low hidden veto transfer & RET (Retention-accountability court) & 0.570 & WFR (Weak-form review) \\
\bottomrule
\end{tabular}
\begin{minipage}{0.96\linewidth}
\footnotesize Notes: Profile scores use alternative display weights. They show that the preferred design can change when the reader prioritizes rights, dialogue, legal stability, low cost, or anti-veto-relocation. These profiles are a first-order weight-sensitivity display; the stress-case sensitivity campaign and full interval files remain in the replication package.
\end{minipage}
\end{table}


Table~\ref{tab:weight-robustness} pushes that test further by sampling random normative weights over twelve baseline outputs. Architecture leadership is unstable across the sampled profiles, which supports the paper's emphasis on tradeoff maps rather than rank ordering. Weak-form review leads most sampled mechanism profiles, but that is not a design recommendation. It benefits from assigned final-veto and response terms. Section~\ref{sec:historical} separately tests legislative uptake, fragmentation and other context assumptions; this new registered exercise does not replace the original baseline comparison or its weight definitions.

% Auto-generated by paper/scripts/generate_tables.py
\begin{table}[hbt!]
\centering
\caption{Random-weight robustness over baseline outputs.}
\label{tab:weight-robustness}
\scriptsize
\setlength{\tabcolsep}{2pt}
\begin{tabular}{@{}L{0.13\linewidth}L{0.22\linewidth}rL{0.20\linewidth}rL{0.17\linewidth}@{}}
\toprule
Design set & Most frequent leader & Share & Next leaders & Mean margin & Reading \\
\midrule
Architectures & RET (Retention-accountability court) & 41.5\% & COUNC 23.5\%; HYB 17.5\%; EMR 17.3\% & 0.004 & Architecture leadership is fragile under broad normative weights. \\
Mechanisms & WFR (Weak-form review) & 92.0\% & RIS 7.9\%; MLR 0.1\% & 0.023 & Weak-form dominance should be tested against legislative-response assumptions. \\
\bottomrule
\end{tabular}
\begin{minipage}{0.96\linewidth}
\footnotesize Notes: The table samples 1000 random normalized weight vectors over twelve outputs: rights protection, responsiveness, legitimacy, legal stability, compliance, inverse conflict, inverse veto-relocation risk, access, implementation capacity, inverse cost, transplant feasibility, and inverse political-culture sensitivity. The margins are differences between the top two designs under each sampled vector.
\end{minipage}
\end{table}


Four patterns emerge from the original experiments. First, emergency-process reforms produce clear direct movement on emergency irregularity, but their substantive interpretation depends on downstream rights-claimant success, legitimacy, compliance, conflict, and cost. Second, appointment, size, and selection reforms mostly act indirectly through composition rather than emergency procedure. Third, cross-checking and council bundles combine legal-performance changes with administrative-load and cost changes; these bundle comparisons do not isolate an upstream displacement effect. Fourth, term limits and appointment systems act through dynamic replacement over time, which is absent when composition is held fixed at initial creation.

The mechanism results add a fifth pattern. Dialogic and front-end bundles can improve the democratic-constitutionalism index through reduced veto-risk terms or modeled legislative responses. Whether legislative uptake explains those advantages requires the response intervention below; simply comparing named bundles cannot answer that question. Rights-impact statements and pre-enactment review lower modeled conflict in some worlds, while their access and screening terms remain assumptions requiring separate scrutiny.

The negative results are equally important. The current-style design remains administratively lean in the model, but its emergency-docket and modeled-legitimacy outputs are weaker in the baseline comparison. Council and cross-checking bundles impose different screening and duplicate-review rules; the magnitude of any resulting burden shift must be measured rather than inferred from those labels. Close display-score differences should not be treated as a fine ranking of institutions.

The value of the comparison is experimental control: each design faces the same generated worlds and dockets. That structure makes it possible to isolate design tradeoffs that would otherwise be obscured by differences in institutional context.

Figure~\ref{fig:cost-score} displays cost against democratic constitutionalism. The current-style court is low cost but has a lower democratic-constitutionalism value in the baseline comparison, while cross-checking, council, and hybrid designs move upward at the price of higher institutional cost. This is the kind of tradeoff that can disappear inside a single aggregate score.

\begin{inlinefigureblock}
% Auto-generated by paper/scripts/generate_figures.py
\begingroup
\setlength{\unitlength}{1mm}
\begin{picture}(128,86)
\scriptsize
\put(33.5,15.0){\color{black!15}\line(0,1){58.0}}
\put(33.5,11.0){\makebox(0,0){0.30}}
\put(56.5,15.0){\color{black!15}\line(0,1){58.0}}
\put(56.5,11.0){\makebox(0,0){0.40}}
\put(79.5,15.0){\color{black!15}\line(0,1){58.0}}
\put(79.5,11.0){\makebox(0,0){0.50}}
\put(102.5,15.0){\color{black!15}\line(0,1){58.0}}
\put(102.5,11.0){\makebox(0,0){0.60}}
\put(22.0,15.0){\color{black!15}\line(1,0){92.0}}
\put(18.0,15.0){\makebox(0,0)[r]{0.55}}
\put(22.0,29.5){\color{black!15}\line(1,0){92.0}}
\put(18.0,29.5){\makebox(0,0)[r]{0.60}}
\put(22.0,44.0){\color{black!15}\line(1,0){92.0}}
\put(18.0,44.0){\makebox(0,0)[r]{0.65}}
\put(22.0,58.5){\color{black!15}\line(1,0){92.0}}
\put(18.0,58.5){\makebox(0,0)[r]{0.70}}
\put(22.0,73.0){\color{black!15}\line(1,0){92.0}}
\put(18.0,73.0){\makebox(0,0)[r]{0.75}}
\put(22.0,15.0){\line(1,0){92.0}}
\put(22.0,15.0){\line(0,1){58.0}}
\put(26,79){\makebox(0,0)[l]{\rule{1.8mm}{1.8mm} court architecture \quad \raisebox{-0.3mm}{\framebox(1.8,1.8){}} mechanism scenario}}
\put(68.0,21.1){\makebox(0,0){\color{black}\rule{2.1mm}{2.1mm}}}
{\color{black!45}\qbezier[8](66.5,22.5)(66.0,24.4)(65.5,26.3)}
\put(64.6,26.3){\makebox(0,0)[r]{CUR*}}
\put(70.3,43.7){\makebox(0,0){\color{black}\rule{1.7mm}{1.7mm}}}
{\color{black!45}\qbezier[8](71.6,45.0)(75.5,50.3)(79.4,55.7)}
\put(80.3,55.7){\makebox(0,0)[l]{EMR}}
\put(100.0,44.0){\makebox(0,0){\color{black}\rule{1.7mm}{1.7mm}}}
\put(103.2,46.0){\makebox(0,0)[l]{DUAL}}
\put(87.1,44.6){\makebox(0,0){\color{black}\rule{1.7mm}{1.7mm}}}
{\color{black!45}\qbezier[8](85.8,45.8)(86.2,48.2)(86.6,50.6)}
\put(86.6,50.6){\makebox(0,0)[c]{COUNC}}
\put(71.7,40.2){\makebox(0,0){\color{black}\rule{1.7mm}{1.7mm}}}
{\color{black!45}\qbezier[8](72.9,39.0)(74.6,36.3)(76.2,33.6)}
\put(77.1,33.6){\makebox(0,0)[l]{OVR}}
\put(89.2,47.8){\makebox(0,0){\color{black}\rule{1.7mm}{1.7mm}}}
{\color{black!45}\qbezier[8](90.4,49.0)(91.3,50.4)(92.3,51.8)}
\put(93.2,51.8){\makebox(0,0)[l]{HYB}}
\put(46.2,60.3){\color{black}\framebox(1.7,1.7){}}
{\color{black!45}\qbezier[8](45.8,62.4)(44.7,63.8)(43.6,65.3)}
\put(42.7,65.3){\makebox(0,0)[r]{WFR}}
\put(51.7,51.9){\color{black}\framebox(1.7,1.7){}}
{\color{black!45}\qbezier[8](51.3,54.0)(50.1,53.8)(48.9,53.7)}
\put(48.0,53.7){\makebox(0,0)[r]{SUS}}
\put(64.2,48.1){\color{black}\framebox(1.7,1.7){}}
{\color{black!45}\qbezier[8](63.8,47.7)(60.8,47.4)(57.9,47.1)}
\put(57.0,47.1){\makebox(0,0)[r]{PRE}}
\put(63.9,47.5){\color{black}\framebox(1.7,1.7){}}
{\color{black!45}\qbezier[8](66.0,49.6)(67.4,50.4)(68.7,51.1)}
\put(69.6,51.1){\makebox(0,0)[l]{ABS}}
\put(82.3,37.9){\color{black}\framebox(1.7,1.7){}}
{\color{black!45}\qbezier[8](84.4,37.5)(85.4,35.7)(86.3,33.8)}
\put(87.2,33.8){\makebox(0,0)[l]{OMB}}
\put(82.8,42.0){\color{black}\framebox(1.7,1.7){}}
{\color{black!45}\qbezier[8](82.4,41.6)(81.5,39.7)(80.5,37.8)}
\put(79.6,37.8){\makebox(0,0)[r]{DEF}}
\put(64.2,55.0){\color{black}\framebox(1.7,1.7){}}
{\color{black!45}\qbezier[8](66.3,57.1)(67.2,61.3)(68.1,65.4)}
\put(69.0,65.4){\makebox(0,0)[l]{RIS}}
\put(54.3,51.9){\color{black}\framebox(1.7,1.7){}}
{\color{black!45}\qbezier[8](53.9,54.0)(54.2,57.8)(54.6,61.7)}
\put(54.6,61.7){\makebox(0,0)[c]{MLR}}
\put(68,4){\makebox(0,0){total institutional cost}}
\put(4,45){\rotatebox{90}{\makebox(0,0){democratic constitutionalism}}}
\end{picture}
\endgroup

\inlinefigurecaption{Baseline institutional cost versus democratic constitutionalism. Filled markers denote court architectures; open markers denote mechanism scenarios. The vertical axis is a composite combining rights, responsiveness, legitimacy, compliance, implementation, access, response credibility, conflict, veto-relocation risk, feasibility, and culture sensitivity. CUR* marks the current-style scenario.}{fig:cost-score}
\end{inlinefigureblock}

Figure~\ref{fig:mechanism-tradeoff} focuses only on mechanism scenarios. The desired region is high democratic constitutionalism and low veto-relocation risk: proposed devices should protect rights and preserve democratic response without hiding final policy control in a less accountable stage.

\begin{figure}[hbt!]
\centering
% Auto-generated by paper/scripts/generate_figures.py
\begingroup
\setlength{\unitlength}{1mm}
\begin{picture}(128,86)
\scriptsize
\put(22.0,15.0){\color{black!15}\line(0,1){58.0}}
\put(22.0,11.0){\makebox(0,0){0.00}}
\put(47.1,15.0){\color{black!15}\line(0,1){58.0}}
\put(47.1,11.0){\makebox(0,0){0.15}}
\put(72.2,15.0){\color{black!15}\line(0,1){58.0}}
\put(72.2,11.0){\makebox(0,0){0.30}}
\put(97.3,15.0){\color{black!15}\line(0,1){58.0}}
\put(97.3,11.0){\makebox(0,0){0.45}}
\put(22.0,15.0){\color{black!15}\line(1,0){92.0}}
\put(18.0,15.0){\makebox(0,0)[r]{0.55}}
\put(22.0,29.5){\color{black!15}\line(1,0){92.0}}
\put(18.0,29.5){\makebox(0,0)[r]{0.60}}
\put(22.0,44.0){\color{black!15}\line(1,0){92.0}}
\put(18.0,44.0){\makebox(0,0)[r]{0.65}}
\put(22.0,58.5){\color{black!15}\line(1,0){92.0}}
\put(18.0,58.5){\makebox(0,0)[r]{0.70}}
\put(22.0,73.0){\color{black!15}\line(1,0){92.0}}
\put(18.0,73.0){\makebox(0,0)[r]{0.75}}
\put(22.0,15.0){\line(1,0){92.0}}
\put(22.0,15.0){\line(0,1){58.0}}
\put(51.6,61.1){\makebox(0,0){\color{black}\rule{1.8mm}{1.8mm}}}
{\color{black!45}\qbezier[8](50.4,62.3)(49.0,64.1)(47.6,65.9)}
\put(47.6,65.9){\makebox(0,0)[r]{WFR}}
\put(56.6,52.7){\makebox(0,0){\color{black}\rule{1.8mm}{1.8mm}}}
{\color{black!45}\qbezier[8](57.8,51.5)(59.2,49.9)(60.6,48.3)}
\put(60.6,48.3){\makebox(0,0)[l]{SUS}}
\put(81.4,45.5){\makebox(0,0){\color{black}\rule{1.8mm}{1.8mm}}}
{\color{black!45}\qbezier[8](82.6,46.7)(84.0,48.4)(85.4,50.1)}
\put(85.4,50.1){\makebox(0,0)[l]{OVC}}
\put(107.5,48.9){\makebox(0,0){\color{black}\rule{1.8mm}{1.8mm}}}
{\color{black!45}\qbezier[8](106.3,47.7)(104.9,45.9)(103.5,44.1)}
\put(103.5,44.1){\makebox(0,0)[r]{PRE}}
\put(74.2,48.4){\makebox(0,0){\color{black}\rule{1.8mm}{1.8mm}}}
{\color{black!45}\qbezier[8](75.4,49.6)(76.9,49.5)(78.4,49.4)}
\put(78.4,49.4){\makebox(0,0)[l]{ABS}}
\put(109.8,38.8){\makebox(0,0){\color{black}\rule{1.8mm}{1.8mm}}}
{\color{black!45}\qbezier[8](111.0,37.6)(112.5,35.8)(114.0,34.0)}
\put(114.0,34.0){\makebox(0,0)[l]{OMB}}
\put(108.0,42.8){\makebox(0,0){\color{black}\rule{1.8mm}{1.8mm}}}
{\color{black!45}\qbezier[8](106.8,41.6)(105.3,39.8)(103.8,38.0)}
\put(103.8,38.0){\makebox(0,0)[r]{DEF}}
\put(106.0,55.9){\makebox(0,0){\color{black}\rule{1.8mm}{1.8mm}}}
{\color{black!45}\qbezier[8](104.8,57.1)(103.3,58.9)(101.8,60.7)}
\put(101.8,60.7){\makebox(0,0)[r]{RIS}}
\put(59.5,52.7){\makebox(0,0){\color{black}\rule{1.8mm}{1.8mm}}}
{\color{black!45}\qbezier[8](60.7,53.9)(62.1,55.7)(63.5,57.5)}
\put(63.5,57.5){\makebox(0,0)[l]{MLR}}
\put(68,4){\makebox(0,0){veto-relocation risk}}
\put(4,45){\rotatebox{90}{\makebox(0,0){democratic constitutionalism}}}
\end{picture}
\endgroup

\caption{Mechanism tradeoff: democratic constitutionalism versus veto-relocation risk.}
\label{fig:mechanism-tradeoff}
\end{figure}

\subsection{Docket-source robustness}

A secondary paired-input comparison imports legislative-output profiles as docket stress signals. Each imported row preserves subject matter when a \texttt{policyDomain} column is available and otherwise infers one from row labels and stress fields. The imported modes change rights threat, public support, legislative mandate, policy-domain mix, capture pressure, volatility, and emergency pressure; they do not change the constitutional-review rules themselves. Because this bridge depends on a separate legislative-simulation data contract, its figure and detailed discussion are reported in the supplement rather than used as a main-text result.

\section{Historical Benchmark and Paired Robustness}
\label{sec:historical}

\subsection{A frozen new-term docket evaluation}

The historical benchmark adds a distinct evidence layer to the original source-range checks. SCDB 2025 Release 01 supplies training records through the 2024 term, while SCDB 2026 Release 01 adds the 2025 term \citep{scdb2025,scdb2026}. The observation unit is one citation-centered dispute, not each consolidated docket, each justice vote, a petition, or only a constitutional merits challenge. The benchmark maps every row into eight mutually exclusive categories, retaining unmapped or residual observations as Other. These new partition definitions do not overwrite the legacy calibration predicates or their denominators.

On September 12, 2026, the benchmark protocol was committed before acquiring new-term case records. Four forecast vectors were then fitted or fixed using the older release and locked before the new-term acquisition. The exposure audit labels 2005--2014 development checks and 2015--2024 backtests as retrospective: old aggregate calibration had already exposed those periods. The 2025 test is new to the recorded pipeline after its forecast freeze, not an assertion of blindness to public events. A separate audit compares all 53 original fields of older records across releases: two historical additions and ten changed matched records are reported separately from the 66 new-term disputes, and none replaces the locked training data.

The empirical candidate uses a 15-term window, a five-term half-life, and a 0.5 pseudocount per category. Its baselines are the frozen generator, expanding historical frequencies, and previous-term frequencies; the latter two also use the fixed pseudocount. No profile uses observed 2025 category shares. Table~\ref{tab:historical-scores} preserves the first test and all three scores. The primary empirical-minus-fixed Brier difference is 0.001704, with a paired conditional 95\% interval of $[-0.007500,0.010336]$. There is no clear conditional advantage. Secondary scores and exposed backtests remain available without substituting a favorable endpoint or refitting after the test. This evaluates docket composition, not merits decisions, constitutional correctness, or the empirical validity of the entire court model.

% Auto-generated by paper/scripts/generate_historical_tables.py
{\small
\setlength{\tabcolsep}{3pt}
\begin{table}[hbt!]
\centering
\caption{Frozen 2025 docket-composition evaluation.}\label{tab:historical-scores}
\begin{tabular}{@{}L{0.36\linewidth}rrr@{}}
\toprule
Forecast & Brier & Log loss & Total variation \\
\midrule
Fixed generator & 0.8088 & 1.8514 & 0.1492 \\
Recency-weighted & 0.8105 & 1.8395 & 0.1391 \\
Expanding history & 0.8119 & 1.9034 & 0.1667 \\
Previous term & 0.8324 & 1.9841 & 0.2335 \\
\bottomrule
\end{tabular}
\par\noindent\begin{minipage}{0.98\linewidth}\footnotesize The 2025 term contains 66 citation-centered disputes. All scores are lower-is-better. Primary recency-minus-fixed Brier difference: 0.001704, paired 95\% interval [-0.007500, 0.010336]. This is within-term, fixed-forecast uncertainty, not between-term uncertainty. The first evaluation is retained; no clear conditional advantage is established.\end{minipage}
\end{table}
}


\subsection{Registered mechanisms, contexts, and estimands}

The institutional study crosses four locked forecast profiles with thirteen contexts and eight designs: current-style review, reasoned emergency review, council review, weak-form review, and four mechanism interventions. Each cell uses 120 runs of 80 cases over four review periods. The full experiment contains 3,993,600 case outcomes and 5,193,352 challenged-object outcomes. Designs within a run face the same complete ordered case/object docket, verified by hashes. Voting streams follow documented design-specific seeds; response interventions use a separately derived seed. Pairing therefore shares docket shocks without claiming identical institutional vote streams.

The interventions change only emergency procedure to the fast-docket setting, council structure to full-court review, or weak-form response to half or zero. The response factor scales credibility before downstream processing and independently thins positive response events. It does not remove the weak-form design's other legal, cost, or score terms. Context sensitivity changes emergency pressure, rights-threat rate, appointment polarization, legislative conflict, implementation capacity, and party fragmentation one at a time, each at two levels around the baseline. This is a fixed sensitivity grid, not a probability distribution over countries or a full interaction design.

The original protocol fixed the grid, component families, interventions, resampling scheme, practical threshold, and preference exercise before the test. The signed component-to-column contract was committed after simulation and the raw audit but before comparative calculations. That later operationalization is disclosed, not represented as independent preregistration. All 19 registered measurements and three additional diagnostics remain visible. Signed base components are stable only if every one of the 52 cell estimates exceeds 0.01 in the expected direction, sign-reversing if both practical signs occur, and mixed/weak otherwise. Unsigned diagnostics receive no directional superiority claim.

Each of the 3,692 reported contrasts averages paired within-run differences. The 2,000-replicate bootstrap resamples whole runs, reusing the same run-index matrix across cells and measures. Its intervals describe conditional Monte Carlo error, not uncertainty about model structure or normative values. The supplement reports all component patterns, all intervention and persistence ranges, and all 71 intervals in the fixed-generator baseline cell. Complete per-cell estimates and standard errors are retained as CSV data. Component intervals are unadjusted descriptive comparisons, not a familywise discovery guarantee.

\subsection{What survives, and what the model cannot establish}

Table~\ref{tab:historical-claims} summarizes the three claim families. Reasoned review lowers modeled emergency irregularity in every cell and raises the total cost index. However, average review time is determined by shared case inputs, so every paired duration difference is exactly zero by construction. The separate delay-cost index cannot rescue a claim about elapsed procedural duration.

% Auto-generated by paper/scripts/generate_historical_tables.py
{\small
\setlength{\tabcolsep}{3pt}
\begin{table}[hbt!]
\centering
\caption{Three institutional claims after the historical robustness study.}\label{tab:historical-claims}
\begin{tabular}{@{}L{0.13\linewidth}L{0.36\linewidth}L{0.43\linewidth}@{}}
\toprule
Claim & Paired difference range & Qualified conclusion \\
\midrule
Emergency & Irregularity: -0.8280 to -0.0456; total cost: 0.0095 to 0.0137. & Lower irregularity; higher cost. Shared review-time input cannot identify procedural duration. Reasons per order require the empty-denominator caveat. \\
Screening & Council upstream cost: 0.0004 to 0.0022; full-court structure intervention: -0.0001 to 0.0001. & Both are below the 0.01 practical threshold throughout. This intervention does not establish substantial upstream displacement. \\
Weak form & Response-off minus full responsiveness: -0.0807 to -0.0189; veto risk: 0.0241 to 0.0872. & Response dependence is partial. Compliance and cost advantages over current persist even with responses disabled. \\
\bottomrule
\end{tabular}
\par\noindent\begin{minipage}{0.98\linewidth}\footnotesize Ranges span 52 fixed forecast/context cells, not confidence intervals. Base contrasts are reasoned or council minus current; ablations are changed design minus its parent. Outputs are model indices/events, not identified reform effects.\end{minipage}
\end{table}
}


Reason-giving also exposes a denominator trap. At low emergency pressure, reasoned review has no emergency order in 116--119 of 120 runs, depending on the forecast profile. The legacy rate assigns zero to those empty runs, producing a numerical sign reversal that mixes order incidence with reasons conditional on an order. The original estimand and classification are retained. Post-estimate pooled counts and rates explain the reversal but do not replace the endpoint or establish a reliable conditional comparison from only one to four orders. Similarly, no legislative replies is not evidence of late or instantaneous replies.

For council review, the upstream-cost difference from current remains below 0.01 in every cell. Switching only council structure to full-court review barely changes that upstream measure. The intervention therefore does not demonstrate substantial upstream displacement, even though direct, capacity, total-cost and administrative-load differences remain visible. Switching from direct-only to full cost accounting changes no practical sign across the 52 cells. These are differently constructed normalized indices, not monetary accounts or an additive omitted-cost estimate.

Turning weak-form response off lowers responsiveness and raises veto risk in all cells. Rights and compliance effects are more context-dependent. The no-response variant nevertheless retains compliance and cost advantages over current, showing that response dependence is partial within this model. Those residual advantages expose other assigned design and outcome terms; they are not empirical evidence that legislative response is unnecessary.

The normative exercise uses the same 1,000 Dirichlet weight vectors on twelve declared measures in every cell, with fixed directions and no fitted rescaling. Full-, half-, and no-response competitions each retain the same three other designs. Weak-form leadership spans 36.1--100\% of weights under full response, 29.9--100\% under half response, and 23.2--99.8\% with responses off. Across base-design pairs, 152 of 312 pair/cell comparisons reverse ordering across sampled preferences. These finite preference shares are not population preferences, welfare estimates, or probabilities that a real institution is best. The shared review-time component contributes no within-cell design discrimination. Thus the study qualifies the original claims rather than selecting a robust institutional winner.

\section{Source-Range Checks, Uncertainty, and Replication}
\label{sec:diagnostics}

The model uses documented external ranges in two ways. Source-range checks require a citable source, nonzero denominator, direct target analogue, and no synthesis or normalized-cost construction. Contextual benchmarks, public-trust ranges, doctrine-mix proxies outside the audited SCDB merits profile, and institutional-cost ranges discipline the comparative setup but are not counted in fit summaries.

Calibration targets identify the court, time period, target measure, target range, observed value when available, sample size, unit, method, reliability, source, source-range flag, and construction note. The current targets include U.S. Supreme Court merits-docket doctrine shares from the Supreme Court Database \citep{scdb2025,scdbIssueCodebook}, emergency-docket context \citep{scotusblogInterim2024}, comparative court intake and throughput proxies, public-trust context \citep{gallupCourtApproval2024}, official-source intake or response rows for Canada, France, the United Kingdom, the European Court of Human Rights, and the CJEU, plus normalized institutional cost profiles. The Canada source rows include a historical Charter-dialogue profile that compares invalidation-conditioned response, override, and reenactment diagnostics to Hogg and Bushell-style response denominators rather than to all simulated cases. The United Kingdom permission-to-appeal row is also mirrored as a case-selection-access proxy; that row disciplines the access measure but is not a direct measure of legal aid, counsel quality, filing cost, standing, or public-interest litigation capacity. Counted rows remain narrow: they require a documented source, denominator, and direct simulator analogue.

The calibration workflow also records promising but not yet source-range-ready findings. Some observations have been incorporated after source verification; others remain excluded because they are secondary-source findings, lack a direct simulator analogue, or need a coding-rule audit. These excluded observations include substituted-law coding for Canadian Charter sequels, conditional French QPC remedy rates, comparative pre-enactment review frequencies, and transplant-feasibility indicators from public-law and governance datasets.

% Auto-generated by paper/scripts/generate_tables.py
\begin{table}[hbt!]
\centering
\caption{Examples of external calibration targets.}
\label{tab:calibration-targets}
\scriptsize
\setlength{\tabcolsep}{3pt}
\begin{tabular}{@{}p{0.26\linewidth}p{0.27\linewidth}rrrl@{}}
\toprule
Court and period & Target & Observed & Range & $n$ & Reliability \\
\midrule
U.S. Supreme Court (2000-2024 terms) & Speech docket share & 0.060 & 0.039--0.080 & -- & high \\
U.S. Supreme Court (2000-2024 terms) & Civil-rights and privacy docket share & 0.164 & 0.137--0.191 & -- & high \\
U.S. Supreme Court (2024-2025 emergency docket) & Substantive emergency application relief rate & 0.440 & 0.310--0.460 & 43 & medium \\
U.S. Supreme Court (2024-2025 emergency docket) & Written explanation share & 0.279 & 0.200--0.360 & 43 & medium \\
U.S. Supreme Court (2024 public-opinion year) & Public court trust and approval & 0.420 & 0.350--0.490 & -- & medium \\
German Federal Constitutional Court (2024) & Constitutional complaint success and admission proxy & 0.009 & 0.006--0.012 & 4640 & medium \\
Supreme Court of Canada (recent annual average) & Leave application grant rate & 0.089 & 0.075--0.105 & 524 & high \\
French Constitutional Council (QPC recent period) & QPC invalidation rate & 0.315 & 0.260--0.360 & 67 & medium \\
Constitutional Court of South Africa (recent annual average) & Petition-to-judgment throughput proxy & 0.141 & 0.110--0.170 & 355 & medium \\
\bottomrule
\end{tabular}
\begin{minipage}{0.96\linewidth}
\footnotesize Notes: The table reports selected calibration targets from the repository's CSV files. Ranges are target intervals, not estimated model effects. Source names, URLs, notes, and validation flags are retained in the replication files.
\end{minipage}
\end{table}


A separate object-disposition diagnostic resolves a Canadian unit mismatch: Table 7 of \citet{morton1992} records 19 nullifications among 50 statute dispositions, not among the study's 100 decisions. Each generated case now contains one or more typed challenge objects, with raw object-level exports. A conditional rule allocates the existing case decision across those objects. Object composition and allocation are unestimated assumptions; matching this source range does not validate the historical docket or independent object-level voting.

The reported confidence intervals reflect run-block variation within the specified model. They do not represent uncertainty about model specification, weights, variable construction, country mapping, or normative assumptions. The bootstrap resamples entire generated-world run blocks rather than isolated cases, which preserves within-run dependence while still using the raw outcome records. Statute-nullification calibration intervals also resample whole run blocks; composition and other calibration intervals retain their documented denominator-based approximations. The alternative normative-profile table, random-weight robustness table, and stress-case sensitivity comparison are therefore more important for interpretation than narrow run-block intervals around one score.

\subsection{External Source-Range Checks}

The source-range comparison is a calibration audit, not empirical authority for the comparative outputs. It runs stylized archetype presets against a shared benchmark docket and compares selected outputs to documented target ranges where a direct analogue exists. Table~\ref{tab:validation-summary} reports the count of mapped targets that fall within range and the largest current miss, if any, for profiles with counted direct targets. A miss may indicate that the simulator abstracts away a case-flow mechanism, that a target is contextual rather than directly comparable, or that the preset should be refined.

% Auto-generated by paper/scripts/generate_tables.py
\begin{table}[hbt!]
\centering
\caption{Validation-style checks against documented target ranges.}
\label{tab:validation-summary}
\scriptsize
\setlength{\tabcolsep}{4pt}
\begin{tabular}{@{}p{0.22\linewidth}p{0.25\linewidth}rrp{0.27\linewidth}@{}}
\toprule
Profile & Scenario preset & In range & Median gap & Largest miss \\
\midrule
U.S. Supreme Court & U.S. Supreme Court benchmark & 1/7 & 0.066 & Emergency-powers merits share \\
U.S. emergency docket & U.S. Supreme Court benchmark & 0/4 & 0.176 & Substantive emergency application relief rate \\
Germany BVerfG & German Federal Constitutional Court & 1/1 & 0.000 & none \\
Canada SCC & Canadian Supreme Court with override context & 1/1 & 0.000 & none \\
France Conseil & French Constitutional Council & 1/1 & 0.000 & none \\
South Africa ConstCourt & South African Constitutional Court & 0/1 & 0.830 & Petition-to-judgment throughput proxy \\
\bottomrule
\end{tabular}
\begin{minipage}{0.96\linewidth}
\footnotesize Notes: The validation campaign runs real-world scenario presets against a shared benchmark docket. The checks count only source-backed target rows marked for validation in the calibration source matrix; synthesis, public-trust, and normalized-cost rows are excluded. They are external diagnostics, not fitted validation estimates.
\end{minipage}
\end{table}


Table~\ref{tab:validation-miss-interpretation} reports the current misses without widening source ranges. Regeneration under a corrected value-based world-seed protocol exposes misses that the earlier process-dependent seed could conceal. Intake-rate misses may reflect denominator conventions, emergency-relief misses procedure-specific applicant mix and merits follow-through, and weak-form response misses the quality and timing of legislative response rather than court merits behavior.

% Auto-generated by paper/scripts/generate_tables.py
\begin{table}[hbt!]
\centering
\caption{Calibration roadmap from the largest source-range misses.}
\label{tab:validation-miss-interpretation}
\scriptsize
\setlength{\tabcolsep}{3pt}
\begin{tabular}{@{}L{0.24\linewidth}L{0.18\linewidth}L{0.21\linewidth}L{0.29\linewidth}@{}}
\toprule
Target & Model vs. source & Category & Interpretation and next step \\
\midrule
Court of Justice of the European Union: Full compliance with rule-of-law rulings & 0.642 vs. 0.534--0.632; gap 0.010 & source-range check & The preset misses this source-range check. Next: Inspect whether the source denominator, simulator metric, and scenario preset are directly comparable. \\
French Constitutional Council: QPC nonconformity rate & 0.299 vs. 0.305--0.324; gap 0.006 & merits-outcome mechanism & The merits-outcome rate is close but outside the source interval for this profile. Next: Retune doctrine severity, referral screening, and invalidation-threshold parameters only after preserving cross-profile comparability. \\
\bottomrule
\end{tabular}
\begin{minipage}{0.96\linewidth}
\footnotesize Notes: The table reports source-range misses when validation-counted checks fall outside documented ranges; when none do, it records the no-current-miss state. The full generated source-range report is retained in \texttt{reports/constitutional-review-validation-v1-misses.csv} and \texttt{.md}.
\end{minipage}
\end{table}


\subsection{Remaining Source-Range Gaps}

The source-range checks are better at finding missing mechanisms than at confirming model fit. A within-range row is a calibration diagnostic rather than country forecasting; out-of-range rows remain explicit empirical-development requirements. Pipeline variables such as generated filing pools, leave filters, lower-court conflict formation, repeat-player selection, and public-interest support are still stylized. Compliance and enforcement measures remain mostly downstream model outputs; the current compliance source check is a narrow CJEU rule-of-law full-compliance diagnostic rather than a whole-court country estimate. Emergency-docket outputs distinguish reasons, vote disclosure, public disagreement, government applicants, merits follow-up, and relief, but they do not yet use a docket-level dataset of applications, responses, referral decisions, and subsequent merits disposition.

The narrowness of Tables~\ref{tab:validation-summary} and \ref{tab:validation-miss-interpretation} is itself a result. Comparative rows are excluded from the source-range count when they rely on synthesis labels, missing denominators, public-trust proxies, unsupported target keys, or normalized-cost assumptions. Empirical extensions should replace each contextual row with a source-specific dataset before using the word validation for country profiles.
Verified provenance is not sufficient by itself. The platform separately flags model-metric gaps when a source reports a reproducible measure whose statistical unit is not exposed by the simulator. Those rows remain outside source-range counts rather than being silently mapped to a broader output, such as treating statute dispositions as all merits decisions.

The supplementary measurement audit adds eight verified, denominator-backed German and South African observations without increasing their validation counts. German complaint success, representation, interim-application demand, and challenged-act composition use distinct procedural units \citep{bverfg2024annual}. South African filing routes contain an unresolved one-petition subtotal discrepancy; dismissals are observed at a July cutoff, and unavailable orders measure source coverage rather than absence of reasons \citep{ally2024docket}. These limitations are retained in the profile cards and replication data.
For that reason, the platform roadmap separates validation-eligible target-family gaps from contextual cost and political-environment source packs. Cost and trust inputs still discipline interpretation, but they should be treated as documented stress context unless future source packs expose raw, reproducible analogues to simulator outputs.

Replication materials separate source code, model inputs, generated outputs, manuscript source, and replication records. The public replication deposit should include source code, model inputs, aggregate outputs, full case-level and interval files or scripts that regenerate them, and instructions for reproducing the simulation outputs.

\section{Implications for Institutional Design}
\label{sec:implications}

The simulation framework points toward a bundle-based account of constitutional review. Independence rules do not have a single effect; they interact with docket control, voting thresholds, review timing, emergency procedure, and enforcement channels. Accountability mechanisms can improve democratic responsiveness but also increase partisan alignment or weaken rights protection if they operate through short renewable terms or retention pressure. Emergency reforms are especially sensitive to whether the reform only demands speed, only demands reasons, or requires a structured merits follow-up pathway.

The model makes administrative accounting assumptions visible, but it does not establish every hypothesized externality. In the registered council comparison, neither the upstream difference nor the isolated structure intervention demonstrates substantial burden displacement. A dual-court bundle's capacity and cost terms likewise should not be read as observed budget or duration effects. Declaration-oriented review changes responsiveness and veto-risk indices, yet turning response off leaves several modeled advantages intact. The practical implication is a requirement for explicit mechanism and measurement tests, not a recommendation to adopt one of these designs.

\section{Limitations and Next Steps}
\label{sec:limits}

The simulator is intentionally abstract. It compresses doctrinal variation, legal culture, party systems, federal design, and executive implementation into variables that can be compared across institutional scenarios. It does not yet model litigant strategy, judicial opinion assignment, actual precedent networks, budget appropriations, media framing, or reform backlash as independently strategic actors. Case-selection access, legislative-response credibility, legislative-response timing, government repeat-player advantage, implementation capacity, legal-transplant feasibility, and political-culture sensitivity are model constructs, not measured country-level estimates. Implementation cost is normalized to the simulator scale and should not be read as a budget forecast. Most comparative rows remain stress-test assumptions until paired with reproducible empirical datasets.

The largest inferential risk is circularity between mechanism definitions and outcome labels. Requiring reasons and connecting relief to merits review directly affect process indices, subject to conditional-denominator effects. Downstream rights, compliance, conflict, legitimacy, response and cost outputs provide additional model diagnostics, but they still depend on assigned coefficients and are not independent causal evidence. The historical study makes three limitations concrete: review time is not design-specific, the council intervention barely moves the hypothesized upstream measure, and weak-form benefits are not fully explained by response. A second risk is parameter arbitrariness: the weights are explicit and sensitivity-tested, not estimated coefficients. A third is comparative overreach. Neither a narrow source-range fit nor a docket-composition test identifies reform effects for named courts.

The most valuable empirical extensions are richer case-flow validation against specific courts and periods; denominator-backed evidence for emergency applications, merits follow-through, legislative response, compliance, enforcement, transplant feasibility, political culture, and case-selection access; and externally sourced indicators for party fragmentation, government control, electoral time pressure, civil-society capacity, public trust, court throughput, and administrative implementation capacity. Future empirical work should replace synthesized targets with denominator-backed datasets, model political-culture variables with external indicators rather than internal assumptions, add explicit implementation budgets for access institutions, and test whether legislative response cycles are credible under different party-system and compliance conditions.

\section*{Supplementary Material}
The supplementary replication package is designed to include the simulator source and tests, calibration and benchmark inputs, source-observation documentation, aggregate outputs, figure- and table-generation scripts, and reproduction commands. Full case-level and interval exports are reproducible from those commands and should be included in the final public deposit or supplied separately if requested.

\section*{Funding Declaration}
No external funding is reported for this manuscript.

\section*{Competing Interests Declaration}
Competing interests: The author declares none.

\section*{Data Availability Statement}
A replication archive will be supplied with the submitted manuscript and, if accepted, deposited in an appropriate public repository. The archive will include source code, model inputs, calibration target-construction inputs, aggregate outputs, manuscript source, and commands needed to reproduce the reported outputs. Full case-level and interval outputs are reproducible from the same commands and should be included in the final public repository or supplied separately if requested during review.

\section*{Ethical Statement}
This project uses synthetic simulation outputs, public court-docket records, and public institutional benchmarks. It does not involve human-subjects research.

\section*{AI Tools Declaration}
OpenAI ChatGPT Deep Research and OpenAI Codex were used from May 1--9, 2026, with further Codex assistance on September 12, 2026. ChatGPT Deep Research was accessed through the ChatGPT web interface at \url{https://chatgpt.com}; the product interface did not expose a more specific model version. It was used for source-discovery prompts, venue-format research, and research-question drafting. OpenAI Codex was accessed through the Codex desktop interface, with the May work identified as GPT-5-family coding assistance. Codex was used for code editing, LaTeX drafting, deterministic figure and table generation, historical-data processing, statistical analysis, replication checks, and manuscript revision. The tools were not trained, fine-tuned, or otherwise modified for this project, and they were not listed as authors. AI outputs were not treated as sources; cited source claims, model choices, figures, code, and manuscript text remain subject to author review and revision.

\begin{thebibliography}{}
% Auto-generated by paper/scripts/generate_tables.py

\bibitem[Baude(2015)]{baude2015}
Baude, William. 2015. ``Foreword: The Supreme Court's Shadow Docket.'' \url{https://chicagounbound.uchicago.edu/journal_articles/8067/}. New York University Journal of Law and Liberty 9: 1.

\bibitem[Bickel(1962)]{bickel1962}
Bickel, Alexander M. 1962. \textit{The Least Dangerous Branch: The Supreme Court at the Bar of Politics}. New Haven, CT: Yale University Press.

\bibitem[Epp(1998)]{epp1998}
Epp, Charles R. 1998. \textit{The Rights Revolution: Lawyers, Activists, and Supreme Courts in Comparative Perspective}. Chicago: University of Chicago Press.

\bibitem[Epstein and Knight(1998)]{epsteinKnight1998}
Epstein, Lee, and Jack Knight. 1998. \textit{The Choices Justices Make}. Washington, DC: CQ Press.

\bibitem[Gallup(2024)]{gallupCourtApproval2024}
Gallup. 2024. ``Approval of U.S. Supreme Court Stalled Near Historical Low.'' \url{https://news.gallup.com/poll/647834/approval-supreme-court-stalled-near-historical-low.aspx}. Accessed May 1, 2026.

\bibitem[Gardbaum(2013)]{gardbaum2013}
Gardbaum, Stephen. 2013. \textit{The New Commonwealth Model of Constitutionalism: Theory and Practice}. Cambridge: Cambridge University Press.

\bibitem[Ginsburg(2003)]{ginsburg2003}
Ginsburg, Tom. 2003. \textit{Judicial Review in New Democracies: Constitutional Courts in Asian Cases}. Cambridge: Cambridge University Press.

\bibitem[Ginsburg and Versteeg(2014)]{ginsburgVersteeg2014}
Ginsburg, Tom, and Mila Versteeg. 2014. ``Why Do Countries Adopt Constitutional Review?.'' \url{https://academic.oup.com/jleo/article/30/3/587/881605}. Journal of Law, Economics, and Organization 30(3): 587--622.

\bibitem[Helmke(2005)]{helmke2005}
Helmke, Gretchen. 2005. \textit{Courts under Constraints: Judges, Generals, and Presidents in Argentina}. Cambridge: Cambridge University Press.

\bibitem[Hirschl(2004)]{hirschl2004}
Hirschl, Ran. 2004. \textit{Towards Juristocracy: The Origins and Consequences of the New Constitutionalism}. Cambridge, MA: Harvard University Press.

\bibitem[Kapiszewski et al.(2013)]{kapiszewski2013}
Kapiszewski, Diana, Gordon Silverstein, and Robert A. Kagan. 2013. \textit{Consequential Courts: Judicial Roles in Global Perspective}. Cambridge: Cambridge University Press.

\bibitem[SCOTUSblog(2025)]{scotusblogInterim2024}
SCOTUSblog. 2025. ``Interim Docket 2024.'' \url{https://www.scotusblog.com/cases/interim-docket/2024/}. Accessed May 1, 2026.

\bibitem[Staton(2010)]{staton2010}
Staton, Jeffrey K. 2010. \textit{Judicial Power and Strategic Communication in Mexico}. Cambridge: Cambridge University Press.

\bibitem[Stone Sweet(2000)]{stoneSweet2000}
Stone Sweet, Alec. 2000. \textit{Governing with Judges: Constitutional Politics in Europe}. Oxford: Oxford University Press.

\bibitem[Supreme Court Database(2025a)]{scdb2025}
Supreme Court Database. 2025a. ``Modern Database: 2025 Release 01.'' \url{https://scdb.la.psu.edu/data/2025-release-01/}. Accessed May 1, 2026.

\bibitem[Supreme Court Database(2025b)]{scdbIssueCodebook}
Supreme Court Database. 2025b. ``Online Codebook: Issue.'' \url{https://scdb.la.psu.edu/online-codebook/issue/}. Accessed May 1, 2026.

\bibitem[Tushnet(1999)]{tushnet1999}
Tushnet, Mark. 1999. \textit{Taking the Constitution Away from the Courts}. Princeton, NJ: Princeton University Press.

\bibitem[Vanberg(2005)]{vanberg2005}
Vanberg, Georg. 2005. \textit{The Politics of Constitutional Review in Germany}. Cambridge: Cambridge University Press.

\bibitem[Vladeck(2023)]{vladeck2023}
Vladeck, Stephen I. 2023. \textit{The Shadow Docket: How the Supreme Court Uses Stealth Rulings to Amass Power and Undermine the Republic}. New York: Basic Books.

\end{thebibliography}

\end{document}
